\chapter{Einleitung}\label{ch:einleitung}

% [Cross-Ref Konsolidierte Doku, K-I.2 2026-05-15]
% Top-Level-Index der Diplomarbeit-Doku: docs/INDEX.md
% Master-Plan Konsolidierungs-Phasen K-A bis K-J: docs/MASTERPLAN_KONSOLIDIERUNG_TERMINE.md
% Begriffsglossar v7 (~161 KANON-Begriffe): docs/glossar/01_begriffsglossar_v7_master.md
% Domaenenmodell v4 (Master): docs/glossar/02_domaenenmodell_v4_master.md
% Termin-Snapshots 1-8: docs/termine_konsolidiert/

\section{Motivation}

Trie-basierte Index-Strukturen sind das Rueckgrat moderner In-Memory-
Datenbanken. Seit dem Adaptive Radix Tree (ART) von Leis et al.
\cite{leis2013art} 2013 ist die Forschung in mehrere Richtungen
ausdifferenziert: Trie-Hoehen-Optimierung (HOT \cite{binna2018hot}),
Hybrid-B+-Strukturen (Masstree \cite{mao2012masstree}, B$^2$-Tree
\cite{schmeisser2022b2tree}), succinct-Repraesentationen (CoCo-Trie
\cite{boffa2024coco}, SuRF \cite{zhang2018surf}) und Self-Tuning
(START \cite{fent2020start}).

Parallel dazu hat sich ein eigener Forschungs-Korpus zu
Cache-Bewusstheit, Prefetching und Allokations-Strategien entwickelt.
Heuristiken wie Cache-Sensitive Search Trees (P11/P12 Rao\&Ross),
Cache-Oblivious Layouts (P16/P17 Bender et al.) und Hierarchical
Prefetching (P27 Zhang et al. ASPLOS 2025) addressieren jeweils
einzelne Achsen der Cache-Pyramide --- aber sind nie systematisch
in einer permutierten Mess-Architektur kombiniert worden.

Diese Arbeit schliesst diese Luecke durch eine drei-schichtige Cache-
Engine-Architektur (CacheEngine $\to$ ExecutionEngine $\to$ SearchEngine)
mit elf Permutations-Achsen, die einen $\sim$5,5 Milliarden grossen
Konfigurations-Raum aufspannen. Als Pruefling fuer die Beitritts-
faehigkeit zum Stand der Technik dient PRT-ART (Probst Redirect
Tree-ART hybrid) mit acht eigenen Bausteine-Schichten.

\section{Forschungsfragen}

\begin{description}
    \item[F1 (Komposition)] Welche elf Algorithmus-Achsen eines
          Cache-bewussten Trie-Index lassen sich kanonisch in einer
          Permutationsmatrix systematisieren, sodass Stand-der-Technik-
          Implementierungen als Profil-Konfigurationen rekonstruiert
          werden koennen?

    \item[F2 (Mess-Methodik)] Wie laesst sich pro Permutation ein fairer,
          reproduzierbarer YCSB-Workload-Lauf durchfuehren, der Profile-
          spezifische Workload-Auswahl (z.\,B. RANGE\_SCAN $\to$ YCSB-E)
          erlaubt und dabei den Mess-Overhead durch ein opt-in
          \texttt{COMDARE\_EXPERIMENT\_MODE}-Flag minimiert?

    \item[F3 (Pruefling-Vergleich)] Wie schneidet PRT-ART als Pruefling
          gegen 8 Tier-1-SOTA (ART, HOT, Masstree, CoCo-Trie, START,
          B$^2$-Tree, Wormhole, SuRF) und 22 Tier-2/3-SOTA (P11--P32)
          unter den drei Pflicht-Messreihen A (Pruefling vs. SOTA),
          B (SOTA-only Vergleichsbasis) und C (Merge-Punkte) ab?
\end{description}

\section{Beitraege}

Die zentralen Beitraege dieser Arbeit sind:

\begin{itemize}
    \item Eine drei-schichtige Architektur (CacheEngine $\to$
          ExecutionEngine $\to$ SearchEngine) mit ABI-stabilen
          C++23-Modul-Interfaces und variadisch-templatierten
          Hybrid-API-Spezialisierungen (1~Param $\Rightarrow$
          \texttt{std::vector}-API; 2~Params $\Rightarrow$
          \texttt{std::map}-API; $N>2 \Rightarrow$
          \texttt{std::map<K, std::tuple<V$_1$\ldots V$_N$>>}).

    \item Eine 11-Achsen-Permutationsmatrix mit
          \texttt{std::variant}-Bausteinen und Compile-Time-Fallback
          (\texttt{resolve\_baustein.hpp}) zwischen Pruefling-Stack und
          Stand-der-Technik-Stack.

    \item 30 SOTA-Algorithmus-Profile als persistierte XML-Konfigurationen
          (8~Tier-1 + 22~Tier-2/3 fuer P11--P32), die der
          \texttt{ExperimentDriver} automatisch zu vorkompilierten
          C++23-Modulen transformiert.

    \item PRT-ART als Pruefling mit acht eigenen Bausteine-Schichten
          (4+2-Pool-Allocator, OLC mit reservierten Bloecken, MultiLevel-
          Layout, PathOriented-Prefetch, DensityTracker, Hypothesis-
          Metriken H1/H2/H3, Inline/External/ChainRef ValueHandle,
          Signaling-Bits-Serialisierung).

    \item Ein Mess-Driver mit Defined-/Full-Mode-Unterscheidung +
          Profile-aware Workload-Routing, der die Habich-Direktive
          ``Messung immer so vollstaendig wie moeglich'' default-haft
          erfuellt.
\end{itemize}

\section{Aufbau der Arbeit}

Kapitel~\ref{ch:sota} fasst den 33-Paper-Stand der Technik zusammen,
gegliedert nach sechs Such-Clustern (A--F) und einem Allokator-Cluster.
Kapitel~\ref{ch:architektur} stellt die Drei-Schichten-Architektur,
das ABI-Interface und die 11-Achsen-Permutationsmatrix vor.
Kapitel~\ref{ch:implementation} dokumentiert die Implementation in den
drei Repositorien (cache-engine, prt-art, Diplomarbeit-Code).
Kapitel~\ref{ch:messverfahren} beschreibt das Mess-Verfahren mit den
drei Pflicht-Messreihen und dem Profile-aware Workload-Routing.
Kapitel~\ref{ch:auswertung} praesentiert die Mess-Resultate und
diskutiert die Hypothesen H1 (Page-Type-Cost), H2 (Code-Qualitaet) und
H3 (Inline-vs-External-Distribution). Kapitel~\ref{ch:fazit} schliesst
mit Fazit und Ausblick.
